\section{Persistent amplification for homogeneous dispersion}
\label{sec:r13-persistent-general}

The exterior commutator estimates yield a stronger statement than one
kinetic excursion: a lower long-time average in actual physical time,
without a global kinetic bound. Let $m>2$, $s=m/2$, and let $a$ satisfy
the positive elliptic, real-even, homogeneous symbol assumptions above.
Write $A(u)=\|a(D)^{1/2}u\|_2^2$, $M(u)=\|u\|_2^2$ and
$P(u)=\operatorname{Im}\int\overline u\nabla u$.

We treat both the local equation $iu_t=a(D)u-|u|^{2m/N}u$ and the
Hartree equation $iu_t=a(D)u-(K*|u|^2)u$ under the preceding Hartree
kernel assumptions: $2<m<N$, a real-even, degree-$-m$ homogeneous
$C^1$ kernel away from zero, positive-definite density form, and
$|K(x)|+|x||\nabla K(x)|\le C|x|^{-m}$. Put
\[
 \kappa=1+m/N\quad\hbox{(local)},\qquad
 \kappa=2\quad\hbox{(Hartree)}.
\]
Their sharp energy inequalities have the common form
\[
 E(u)=A(u)/2-\mathcal F(u)/(2\kappa),\qquad
 \mathcal F(u)\le\kappa(M(u)/M_c)^{\kappa-1}A(u).
\]
Here $\mathcal F=\|u\|_{2\kappa}^{2\kappa}$ locally and is the Hartree
interaction otherwise.

\begin{theorem}[Lower long-time kinetic average]\label{thm:r13-general-average}
For each fixed equation above there are $d_0,c>0$ such that every
forward-global conservative strong $H^s$ solution, for which the
localized identities pass by strong energy-space approximation, with
\[
 P=0,\qquad E=-e<0,\qquad
 0<d=(M/M_c)^{\kappa-1}-1<d_0
\]
satisfies
\begin{equation}\label{eq:r13-general-average}
 \boxed{\quad
 \liminf_{T\to\infty}\frac1T\int_0^T A(u(t))\,dt
 \ge c e d^{-\alpha_m},\qquad
 \alpha_m=1+\frac{m}{2(m-1)}.
 \quad}
\end{equation}
No uniform kinetic bound, radial symmetry, finite spatial moment,
convexity of $a$, or real-sector linearized nondegeneracy is assumed.
The theorem specifies its solution class; it does not assert local
well-posedness for every symbol and dimension.
\end{theorem}

\begin{proof}
Write $A(t)=A(u(t))$, $\delta=M-M_c$, $\gamma=(m-1)/m$.
Sharp interpolation gives $A(t)\ge a_*:=2e/d>0$.
Real polarization supplies, separately at each time, a phase with
$u=e^{i\theta}(X+iY)$, $X,Y$ real, and
$M(Y)\le CdM$, $A(Y)\le CdA(t)$.
The nonlinear potential is real in both models and cancels from the
localized mass derivative. Apply the normalized Fourier mass-commutator
bound after scaling $A$ to one. For $R A(t)^{1/m}\ge1$ it gives
\begin{equation}\label{eq:r13g-variable-density}
 \left|\frac d{dt}\int\chi((x-z)/R)|u|^2\right|
 \le C\sqrt d R^{-1}A(t)^\gamma.
\end{equation}
No time derivative of the chosen phase is taken.

\medskip\noindent\emph{Center and uniform mass tail.}
Fix $\varepsilon_0=M_c/32$ and take $d_0$ small. Normalized
near-optimizer compactness gives a fixed $\overline R$ and centers
$c(t)$ with
\[
 \int_{|x-c(t)|\le\overline R A(t)^{-1/m}}|u(t,x)|^2\,dx
 \ge M_c-\varepsilon_0.
\]
Choose $R_0\ge1$ containing these balls and satisfying
$R_0\ge a_*^{-1/m}$. Introduce $s(t)=\int_0^tA(\tau)^\gamma d\tau$.
Two centers separated by more than $4R_0$ force a cutoff mass change
of at least $M_c-3\varepsilon_0$. Equation
\eqref{eq:r13g-variable-density} excludes this when
$|s(t)-s(t')|\le\tau R_0/\sqrt d$, with fixed small $\tau>0$.
Interpolating sampled centers at these intrinsic-time intervals gives
a locally absolutely continuous $b(t)$ with
\begin{equation}\label{eq:r13g-variable-center}
 |b(t)-c(t)|\le8R_0,\qquad
 |b'(t)|\le C\sqrt d A(t)^\gamma\quad\hbox{a.e.}
\end{equation}
The speed constant is independent of $R_0$ and $e/d$.

For every $\eta>0$ there is $R_\eta<\infty$ such that
\begin{equation}\label{eq:r13g-variable-tail}
 \sup_{t\ge0}\int_{|x-b(t)|>R_\eta}|u(t,x)|^2\,dx
 \le\delta+\eta.
\end{equation}
Here is the compactness argument without a kinetic upper bound.
Suppose bad times $t_k$ and radii $R_k\to\infty$ exist. The bad mass
lies also outside $B(c(t_k),R_k-8R_0)$. Set $A_k=A(t_k)$,
$\ell_k=A_k^{-1/m}$ and
$v_k(y)=\ell_k^{N/2}u(t_k,c(t_k)+\ell_k y)$.
Then $A(v_k)=1$, $M(v_k)=M$, $E(v_k)=-e/A_k\in[-d/2,0)$,
and a fixed ball contains $M_c-\varepsilon_0$ mass.
A weak $H^s$ limit $v$, with strong local $L^2$ convergence, therefore
has $M(v)\ge M_c-\varepsilon_0>0$. The normalized bad radius tends
to infinity, since
$(R_k-8R_0)\ell_k^{-1}\ge(R_k-8R_0)a_*^{1/m}\to\infty$.
Consequently $M(v)\le M_c-\eta$.
For $w_k=v_k-v$, mass and energy split:
\[
 E(v_k)=E(v)+E(w_k)+o(1),\qquad
 M(w_k)\le\delta+\varepsilon_0+o(1)<M_c.
\]
Locally the nonlinear splitting is Brezis--Lieb. For Hartree, the
Riesz bound and HLS imply that $v\overline{w_k}$ vanishes in the
density space, while $|w_k|^2$ converges weakly to zero there. First
approximate $v$ by bounded compactly supported functions and use local
compactness; then use uniform Sobolev and HLS bounds. This removes all
mixed interaction terms. Sharp interpolation gives $E(w_k)\ge0$
eventually and $E(v)>0$, because $v\ne0$ is strictly subcritical,
contradicting $E(v_k)<0$. This proves \eqref{eq:r13g-variable-tail}.

\medskip\noindent\emph{Variable-scale virial.}
Let $g=a(D)^{1/2}u$ and use the compact vector field $f_R$ of the
exterior commutator lemma. Define
\[
 J_R=2\operatorname{Im}\int f_R(x-b(t))\cdot\overline u\nabla u,
 \quad K_R=\int_{|x-b|>8R}|g|^2,
 \quad H_R=\int_{R/4<|x-b|<8R}|g|^2.
\]
Scaling the normalized exterior commutator and cutoff-energy estimates
back to $A(t)$ gives errors $CR^{-1}A^\gamma$ and
$CR^{-2}A^{(m-2)/m}$. Write
\[
 q_R(t)=R^{-1}A(t)^\gamma+R^{-2}A(t)^{(m-2)/m}+R^{-m}.
\]
The last term is needed only for separated Hartree interactions.
For $R\ge4R_\eta$, $R\ge a_*^{-1/m}$ and $\mu=\delta+\eta$,
the preceding localized identities give
\begin{align*}
 J_R'\le{}&-4me-2mK_R+CH_R+Cq_R\\
 &+C\mu^{\kappa-1}(K_R+H_R+q_R)\\
 &+C|b'|\mu^{1-1/m}(K_R+H_R+q_R)^{1/m}.
\end{align*}
The constant center-motion term is $-2b'\cdot P=0$. Its remainder is
bounded using $\|\zeta_Ru\|_2\|\nabla(\zeta_Ru)\|_2$ and Fourier
interpolation, since $s>1$. For Hartree, separated inner--outer
interactions are $O(R^{-m}M^2)$; the remaining exterior interaction
is bounded by $C\mu A(\zeta_Ru)$.

Take $d_0,\eta$ small enough to absorb the nonlinear $K_R$ term.
Young's inequality absorbs the current's $K_R$ term and bounds the
other parts by $C(H_R+q_R)+C\mu|b'|^{m/(m-1)}$.
Equation \eqref{eq:r13g-variable-center} implies
\begin{equation}\label{eq:r13g-average-virial}
 J_R'\le-4me+CH_R+C\mu d^{m/[2(m-1)]}A(t)+Cq_R(t).
\end{equation}
The endpoint estimate is $|J_R(t)|\le CR A(t)^{1/m}$.

\medskip\noindent\emph{Physical-time averaging.}
Let $L$ be the left side of \eqref{eq:r13-general-average}; only finite
$L$ needs proof. Choose $T_n\to\infty$ with averages tending to $L$.
There is $t_n\in[T_n-\sqrt{T_n},T_n]$ with
$A(t_n)\le C(L+1)\sqrt{T_n}$. Therefore $t_n/T_n\to1$, the
averages through $t_n$ tend to $L$, and $A(t_n)^{1/m}/t_n\to0$.
For fixed $K<\infty$ and large $R_*$, the annuli at radii
$R_j=64^jR_*$, $0\le j<K$, are disjoint, so
$\sum_jH_{R_j}(t)\le A(t)$.
For each $n$ choose one radius with average annular energy at most
$K^{-1}$ times the average of $A$ and integrate
\eqref{eq:r13g-average-virial}.
The chosen radii stay bounded for fixed $K,R_*$, so endpoints divided
by $t_n$ vanish. Jensen's inequality bounds the fractional powers of
$A$ in $q_R$. The time limit yields
\[
 4me\le CL/K+C\mu d^{m/[2(m-1)]}L
       +CR_*^{-1}L^\gamma+CR_*^{-2}L^{(m-2)/m}+CR_*^{-m}.
\]
Let $K\to\infty$, then $R_*\to\infty$, then $\eta\downarrow0$.
Since $\delta=M_c[(1+d)^{1/(\kappa-1)}-1]\le Cd$, this proves
\eqref{eq:r13-general-average}.
\end{proof}

\begin{corollary}[Repeated amplification and standing-wave instability]
\label{cor:r13-general-repeat}
For any real sharp standing wave $Q$ with $M(Q)=M_c$, take
$u_d(0)=(1+d)^{1/[2(\kappa-1)]}Q$.
Every global solution in the stated class satisfies
\[
 \liminf_{T\to\infty}\frac1T\int_0^T
  \frac{A(u_d(t))}{A(u_d(0))}\,dt
 \ge c d^{-m/[2(m-1)]}.
\]
Every prescribed finite amplification factor is thus exceeded at
arbitrarily late times for small enough $d$. Under the usual local
continuation alternative, each real sharp standing wave is nonlinearly
orbitally unstable modulo phase and translation.
\end{corollary}
\begin{proof}
The exact amplitude identity gives $e=dA(u_d(0))/2$.
The initial perturbations converge to $Q$ in $H^s$, whereas a large
kinetic excursion leaves a fixed neighborhood of its phase--translation
orbit. Finite-time breakdown also forces such an excursion under the
continuation alternative.
\end{proof}

If a particular orbit has $B=\sup_tA(t)<\infty$ and
$L_*=ce d^{-\alpha_m}$, then the set $\{t:A(t)\ge L_*/2\}$ has
lower asymptotic density at least $L_*/(2B)$. Without this extra
boundedness, the theorem still permits sparse high spikes. For $m=4$,
$\alpha_m=5/3$ and the perturbation exponent is $2/3$. The result does
not by itself exclude every bounded global negative-energy orbit.
